% Part: model-theory
% Chapter: lindstrom
% Section: lindstrom-proof

\documentclass[../../../include/open-logic-section]{subfiles}

\begin{document}

\olfileid{mod}{lin}{prf}

\olsection{مبرهنة ليندستروم}

\begin{lem}
\ollabel{lem:lindstrom}
لتكن $!E \in L(\Lang{L})$، ولتكن $\Lang{L}$ منتهية. فإن
وجد $n \in \Nat$ تتحقق معه، لكل !!{structure}s اثنتين $\Struct{M}$ و
$\Struct{N}$، القضية الآتية: إذا كان $\Struct{M} \equiv_n \Struct{N}$ و
$\Struct{M} \models_L !E$، فإن $\Struct{N} \models_L !E$ أيضًا؛
فإن $!E$ تكافئ~!!{sentence} من الرتبة الأولى؛ أي توجد $!D$ من الرتبة
الأولى بحيث $\Mod(L){!E}=\Mod(L){!D}$.
\end{lem}

\begin{proof}
ليكن $n$ بحيث تتفق أي !!{structure}s متكافئتين حتى~$n$، $\Struct{M}$
و$\Struct{N}$، على القيمة المسندة إلى~$!E$. وقد ثبت في
\olref[bas][pis]{prop:qr-finite} أن !!{language} منتهية لا تشتمل، إذا سوّينا بين المتكافئات
منطقيًا، إلا على عدد منتهٍ من !!{sentence}s الرتبة الأولى التي لا
تتجاوز رتبة كمها~$n$. ولكل !!{structure} ثابتة~$\Struct{M}$، لتكن
$!D_{\Struct{M}}$ اقتران ممثل واحد من كل فئة تكافؤ منطقي للـ!!{sentence}s
$!A$ من الرتبة الأولى الصادقة في~$\Struct{M}$ والتي تحقق
$\QuantRank{!A} \le n$ (فهذا الاقتران منتهٍ)، بحيث
$\Struct{N} \models !D_{\Struct{M}}$ إذا وفقط إذا كان
$\Struct{N} \equiv_n \Struct{M}$. ثم ضع
$!D = \textstyle\bigvee
\Setabs{!D_{\Struct{M}}}{\Struct{M} \models_L !E}$؛ وهذا الانفصال منتهٍ
أيضًا حتى التكافؤ المنطقي.

فنحصل على $\Mod(L){!E}=\Mod(L){!D}$، وبيان الوجهين كما يأتي. إن كان
$\Struct{N} \models_L !D$، وجد $\Struct{M} \models_L !E$ بحيث
$\Struct{N} \models !D_{\Struct{M}}$، ومن ثم
$\Struct{N} \models_L !E$ أيضًا (بفرض اللمة). وأما العكس، فإن كان
$\Struct{N} \models_L !E$، كان $!D_\Struct{N}$ أحد حدود الانفصال في
$!D$، ولما كان $\Struct{N} \models !D_\Struct{N}$، كان كذلك
$\Struct{N} \models_L !D$.
\end{proof}

\begin{thm}[مبرهنة ليندستروم]
  \ollabel{thm:lindstrom} إذا كان $\tuple{L,\models_L}$ منطقًا مجردًا سويًا
  يجمع خاصيتي التراص ولوفنهايم--سكولم، فإن
  $\tuple{L,\models_L} \le \tuple{F,\models}$ (ومن ثم يكون
  $\tuple{L,\models_L}$ مكافئًا لمنطق الرتبة الأولى).
\end{thm}

\begin{proof}
ترجع المسألة، بمقتضى \olref{lem:lindstrom}، إلى إثبات أنه لكل
$!E \in L(\Lang{L})$، حيث $\Lang{L}$ منتهية، يوجد $n \in \Nat$ بحيث
إنه لأي !!{structure}s اثنتين $\Struct{M}$ و$\Struct{N}$: إذا كان
$\Struct{M} \equiv_n \Struct{N}$، اتفقت $\Struct{M}$ و$\Struct{N}$
على~$!E$. فعندئذ تكافئ $!E$~!!{sentence} من الرتبة الأولى، فيترتب على ذلك
$\tuple{L,\models_L} \le \tuple{F,\models}$. وحيث إن البحث في
!!{language} منتهية علائقية محضة، تجيز لنا
\olref[bas][pis]{thm:b-n-f} إحلال القضية الجبرية المقابلة
$I_n(\emptyset,\emptyset)$ محل القضية
$\Struct{M} \equiv_n \Struct{N}$.

فلنأخذ $!E$، ولنفرض خلاف المطلوب: أنه لكل $n$ توجد
!!{structure}s~$\Struct{M}_n$ و$\Struct{N}_n$ بحيث
$I_n(\emptyset,\emptyset)$، لكن، مثلًا،
$\Struct{M}_n \models_L !E$ و$\Struct{N}_n \not\models_L !E$. وتجيز لنا خاصية
التشاكل أن نفرض أن جميع الـ$\Struct{M}_n$ تفسر ثوابت اللغة
بالأشياء نفسها؛ ثم إن اللغة لا تحتوي إلا عددًا منتهيًا من
الـ!!{sentence}s الذرية؛ فلنا أن نفترض اتفاق هذه البنى في الـ!!{sentence}s الذرية
التي تحققها، وإلا اقتصرنا على متتالية جزئية من الـ$\Struct{M}$. ولتكن
$\Struct{M_\infty}$ اتحاد جميع الـ$\Struct{M}_n$، أي الـ!!{structure}
الصغرى الوحيدة التي تحتوي كل $\Struct{M}_n$ بنية جزئية. وكما في برهان
\olref[lsp]{thm:abstract-p-isom}، لتكن $\Struct{M_\infty}^*$ توسيع
$\Struct{M_\infty}$ ذي الـ!!{domain}
$\Domain{M_\infty} \cup \Domain{M_\infty}^{<\omega}$، في
الـ!!{language} الموسعة التي تضم محمولي الوصل $P$ و$Q$.

ونجري البناء نفسه لتعريف $\Struct{N}_n$ و$\Struct{N_\infty}$ و
$\Struct{N_\infty}^*$. ولتكن الآن $\Struct{U}$ هي الـ!!{structure} التي
يضم !!{domain}ها !!{domain}s~$\Struct{M_\infty}^*$ و
$\Struct{N_\infty}^*$، وكذلك الأعداد الطبيعية~$\Nat$ مع ترتيبها
الطبيعي~$\le$، في الـ!!{language} ذات محمولات إضافية تمثل الـ!!{domain}s
$\Domain{M_\infty}$ و$\Domain{N_\infty}$ و
$\Domain{M_\infty}^{<\omega}$ و$\Domain{N_\infty}^{<\omega}$، فضلًا
عن محمولات ترمّز مجالات $\Struct{M}_n$ و$\Struct{N}_n$ بمعنى أن:
\begin{align*}
  \Domain{M_n} & = \Setabs{a \in \Domain{M_\infty}}{R(a, n)}; &
  \Domain{N_n} & = \Setabs{a \in \Domain{N_\infty}}{S(a,n)}; \\
  \Domain{M_\infty}^{<\omega}_n & = \Setabs{a \in \Domain{M_\infty}^{<\omega}}{R(a,n)}; &
  \Domain{N_\infty}^{<\omega}_n & = \Setabs{a \in \Domain{N_\infty}^{<\omega}}{S(a,n)}.
\end{align*}
وللـ!!{structure}~$\Struct{U}$ أيضًا علاقة ثلاثية~$J$ بحيث
$J(n,\mathbf{a},\mathbf{b})$ إذا وفقط إذا كان
$I_n(\mathbf{a},\mathbf{b})$.

ويمكن الآن صوغ !!a{sentence}~$!D$ من المنطق~$L$، في الـ!!{language}
$\Lang{L}$ الموسعة بـ$R$ و$S$ و$J$ وغيرها، تقول إن $\le$ ترتيب خطي
منفصل له عنصر أول ولا عنصر أخير، وإن
$\Struct{M}_n \models_L !E$ و$\Struct{N}_n \not\models_L !E$، وإنه لكل
$n$ في الترتيب يكون $J(n,\mathbf{a},\mathbf{b})$ إذا وفقط إذا كان
$I_n(\mathbf{a},\mathbf{b})$.

نلحق باللغة ثابتًا جديدًا~$n^*$، ثم نضم إلى $!D$ نمطًا يقتضي أن
$\bar n < n^*$ لكل $n \in \Nat$ قياسي (وبصورة مكافئة، يفرض لكل عدد
قياسي~$n$ أن تحت $n^*$ ما لا يقل عن $n$ عنصرًا). ولا يعسر تحقيق أي جزء منتهٍ من هذا
النمط مع $!D$ في $\Struct{U}$؛ إذ يكفي اختيار عنصر قياسي بالغ الكبر.
فتقتضي خاصية التراص وجود نموذج~$\Struct{V}$ لـ$!D$ وللنمط بأسره، وفي
ترتيبه عنصر~$n^*$ يعلو كل عنصر قياسي. وبوجه خاص، تحتوي
$\Struct{V}$~!!{structure}s جزئيتين $\Struct{M_{n^*}}$ و
$\Struct{N_{n^*}}$ بحيث $\Struct{M_{n^*}} \models_L !E$ و
$\Struct{N_{n^*}} \not\models_L !E$. وحينئذ نعرّف مجموعة
$\mathcal{I}$ من أزواج متتاليات طولها~$k$ من $\Domain{M_{n^*}}$ و
$\Domain{N_{n^*}}$، بوضع
$\tuple{\mathbf{a},\mathbf{b}} \in \mathcal{I}$ إذا وفقط إذا كان
$J(n^*-k,\mathbf{a},\mathbf{b})$، حيث $k$ طول $\mathbf{a}$ و
$\mathbf{b}$. وحيث إن $n^*$ يعلو كل عدد قياسي، لزم أن
$n^*-k>0$ لكل $k$ قياسي، وتشهد المجموعة $\mathcal{I}$ أن
$\Struct{M_{n^*}} \simeq_p \Struct{N_{n^*}}$. غير أن
\olref[lsp]{thm:abstract-p-isom} تقضي بأن $\Struct{M_{n^*}}$ مكافئة لـ
$\Struct{N_{n^*}}$ في~$L$؛ فتناقض اختلافهما المفروض، وثبت المطلوب.
\end{proof}

\end{document}
